% ======================================================================
%  preamble.tex --- 本書共通のプリアンブル（全章がこれを共有する）
%  章ファイルの中でパッケージを読み込んだりマクロを再定義したりしないこと。
% ======================================================================
\usepackage[haranoaji,deluxe]{luatexja-preset}
\usepackage{amsmath,amssymb,amsthm,mathtools,bm}
\usepackage{graphicx}
\usepackage[dvipsnames,svgnames]{xcolor}
\usepackage{booktabs,multirow,array,tabularx,makecell}
\usepackage{enumitem}
\usepackage{caption,subcaption}
\usepackage{float}
\usepackage{tikz}
\usetikzlibrary{arrows.meta,positioning,calc,fit,backgrounds,matrix,
  shapes.geometric,shapes.misc,shapes.multipart,decorations.pathreplacing,
  decorations.pathmorphing,patterns,chains,shadows}
\usepackage{pgfplots}
\pgfplotsset{compat=1.18}
\usepgfplotslibrary{fillbetween}
\usepackage{listings}
\usepackage[most]{tcolorbox}
\usepackage{fancyhdr}
\usepackage{makeidx}
\usepackage[unicode,hidelinks,bookmarksnumbered,pdfusetitle]{hyperref}
\usepackage{url}

% ---------------------------------------------------------------- 色
\definecolor{bookblue}{HTML}{1F4E79}
\definecolor{bookaccent}{HTML}{C0504D}
\definecolor{bookgreen}{HTML}{2E7D32}
\definecolor{bookorange}{HTML}{E08A00}
\definecolor{bookpurple}{HTML}{6A4C93}
\definecolor{bookgray}{HTML}{5F6B7A}
\definecolor{codebg}{HTML}{F6F8FA}
\definecolor{outbg}{HTML}{FBFAF3}
\colorlet{figblue}{bookblue!18}
\colorlet{figred}{bookaccent!18}
\colorlet{figgreen}{bookgreen!18}
\colorlet{figorange}{bookorange!22}
\colorlet{figpurple}{bookpurple!18}
\colorlet{figgray}{black!8}

% ---------------------------------------------------------------- 版面・見出し
\setlength{\headheight}{14pt}
\pagestyle{fancy}
\fancyhf{}
\fancyhead[L]{\small\sffamily\color{bookgray}\leftmark}
\fancyhead[R]{\small\sffamily\color{bookgray}\thepage}
\renewcommand{\headrulewidth}{0.3pt}
\fancypagestyle{plain}{\fancyhf{}\fancyfoot[C]{\small\sffamily\thepage}%
  \renewcommand{\headrulewidth}{0pt}}
% 第 0 章でも「図 0.1」のように章番号を付ける。柱は「第 3 章　タイトル」。
\renewcommand{\thefigure}{\thechapter.\arabic{figure}}
\renewcommand{\thetable}{\thechapter.\arabic{table}}
\renewcommand{\theequation}{\thechapter.\arabic{equation}}
\makeatletter
\renewcommand{\chaptermark}[1]{\markboth{\@chapapp\thechapter\@chappos\quad #1}{}}
\makeatother
\setcounter{tocdepth}{1}
\setcounter{secnumdepth}{2}
\captionsetup{font=small,labelfont={bf,sf},labelsep=quad}
\setlist{itemsep=0.2ex,topsep=0.6ex,parsep=0.2ex}
\allowdisplaybreaks[2]

% ---------------------------------------------------------------- 数式マクロ
% ベクトルは太字小文字、行列は太字大文字。
\newcommand{\vect}[1]{\bm{#1}}
\newcommand{\mat}[1]{\bm{#1}}
\newcommand{\va}{\bm{a}}\newcommand{\vb}{\bm{b}}\newcommand{\vc}{\bm{c}}
\newcommand{\ve}{\bm{e}}\newcommand{\vg}{\bm{g}}\newcommand{\vh}{\bm{h}}
\newcommand{\vk}{\bm{k}}\newcommand{\vm}{\bm{m}}\newcommand{\vo}{\bm{o}}
\newcommand{\vp}{\bm{p}}\newcommand{\vq}{\bm{q}}\newcommand{\vr}{\bm{r}}
\newcommand{\vs}{\bm{s}}\newcommand{\vu}{\bm{u}}\newcommand{\vv}{\bm{v}}
\newcommand{\vw}{\bm{w}}\newcommand{\vx}{\bm{x}}\newcommand{\vy}{\bm{y}}
\newcommand{\vz}{\bm{z}}\newcommand{\vzero}{\bm{0}}\newcommand{\vone}{\bm{1}}
\newcommand{\vtheta}{\bm{\theta}}\newcommand{\vphi}{\bm{\phi}}
\newcommand{\vdelta}{\bm{\delta}}\newcommand{\vmu}{\bm{\mu}}
\providecommand{\vgamma}{\bm{\gamma}}\providecommand{\vbeta}{\bm{\beta}}
\newcommand{\mA}{\bm{A}}\newcommand{\mB}{\bm{B}}\newcommand{\mC}{\bm{C}}
\newcommand{\mD}{\bm{D}}\newcommand{\mE}{\bm{E}}\newcommand{\mH}{\bm{H}}
\newcommand{\mI}{\bm{I}}\newcommand{\mJ}{\bm{J}}\newcommand{\mK}{\bm{K}}
\newcommand{\mM}{\bm{M}}\newcommand{\mP}{\bm{P}}\newcommand{\mQ}{\bm{Q}}
\newcommand{\mR}{\bm{R}}\newcommand{\mS}{\bm{S}}\newcommand{\mU}{\bm{U}}
\newcommand{\mV}{\bm{V}}\newcommand{\mW}{\bm{W}}\newcommand{\mX}{\bm{X}}
\newcommand{\mY}{\bm{Y}}\newcommand{\mZ}{\bm{Z}}
\newcommand{\R}{\mathbb{R}}                    % 実数全体
\newcommand{\N}{\mathbb{N}}                    % 自然数全体
\newcommand{\E}{\mathbb{E}}                    % 期待値
\newcommand{\Var}{\operatorname{Var}}          % 分散
\newcommand{\calD}{\mathcal{D}}                % データセット
\newcommand{\calL}{\mathcal{L}}                % 損失関数
\newcommand{\calV}{\mathcal{V}}                % 語彙
\newcommand{\calS}{\mathcal{S}}                % 状態集合
\newcommand{\calA}{\mathcal{A}}                % 行動集合
\newcommand{\calN}{\mathcal{N}}                % 正規分布
\newcommand{\calO}{\mathcal{O}}                % オーダー記法
\newcommand{\T}{\top}                          % 転置
\newcommand{\pdata}{p_{\mathrm{data}}}         % データの真の分布
\newcommand{\ptheta}{p_{\theta}}               % モデルの分布
\newcommand{\pitheta}{\pi_{\theta}}            % 方策
\newcommand{\KL}[2]{D_{\mathrm{KL}}\!\left(#1\,\middle\|\,#2\right)}
\newcommand{\pd}[2]{\frac{\partial #1}{\partial #2}}   % 偏微分
\newcommand{\dd}[2]{\frac{d #1}{d #2}}                 % 常微分
\newcommand{\softmax}{\operatorname{softmax}}
\newcommand{\relu}{\operatorname{ReLU}}
\newcommand{\gelu}{\operatorname{GELU}}
\newcommand{\sigmoid}{\sigma}
\DeclareMathOperator*{\argmax}{arg\,max}
\DeclareMathOperator*{\argmin}{arg\,min}
\newcommand{\norm}[1]{\left\lVert #1 \right\rVert}
\newcommand{\abs}[1]{\left\lvert #1 \right\rvert}
\newcommand{\inner}[2]{\left\langle #1, #2 \right\rangle}
\newcommand{\defeq}{\coloneqq}                 % 「左辺を右辺で定義する」
\newcommand{\dmodel}{d}                        % 埋め込み次元
\newcommand{\dhead}{d_{\mathrm{head}}}
\newcommand{\dff}{d_{\mathrm{ff}}}

% ---------------------------------------------------------------- 用語・参照
% \term{日本語}{English}: 初出の専門用語。太字にして索引にも載せる。
\newcommand{\term}[2]{\textbf{#1}（#2）\index{#1@#1（#2）}}
% \jterm{日本語}: 英語を添えない初出用語。
\newcommand{\jterm}[1]{\textbf{#1}\index{#1}}
\newcommand{\chapref}[1]{第\ref{#1}章}
\newcommand{\secref}[1]{\ref{#1}節}
\newcommand{\figref}[1]{図\ref{#1}}
\newcommand{\tabref}[1]{表\ref{#1}}
\newcommand{\coderef}[1]{コード\ref{#1}}
\newcommand{\appref}[1]{付録\ref{#1}}
\newcommand{\partref}[1]{第\ref{#1}部}
\newcommand{\code}[1]{\texttt{\detokenize{#1}}}   % 文中のコード片
\newcommand{\file}[1]{\texttt{\detokenize{#1}}}   % ファイル名

% ---------------------------------------------------------------- 定理環境
\theoremstyle{definition}
\newtheorem{definitioninner}{定義}[chapter]
\newtheorem{theoreminner}{定理}[chapter]
\newtheorem{exampleinner}{例}[chapter]
\renewcommand{\proofname}{\textbf{証明}}

\tcbset{bookbox/.style={enhanced,breakable,sharp corners,boxrule=0pt,
  left=3mm,right=3mm,top=1.5mm,bottom=1.5mm,before skip=3mm,after skip=3mm}}

% 定義: \begin{definition}[名前] ... \end{definition}
\newenvironment{definition}[1][]{%
  \begin{tcolorbox}[bookbox,colback=figblue!45,borderline west={2.5pt}{0pt}{bookblue}]
  \ifx\relax#1\relax\begin{definitioninner}\else\begin{definitioninner}[#1]\fi}%
  {\end{definitioninner}\end{tcolorbox}}
% 定理: \begin{theorem}[名前] ... \end{theorem}
\newenvironment{theorem}[1][]{%
  \begin{tcolorbox}[bookbox,colback=figpurple!45,borderline west={2.5pt}{0pt}{bookpurple}]
  \ifx\relax#1\relax\begin{theoreminner}\else\begin{theoreminner}[#1]\fi}%
  {\end{theoreminner}\end{tcolorbox}}
% 例: \begin{example}[名前] ... \end{example}
\newenvironment{example}[1][]{%
  \ifx\relax#1\relax\begin{exampleinner}\else\begin{exampleinner}[#1]\fi}%
  {\hfill$\triangleleft$\end{exampleinner}}

% ---------------------------------------------------------------- 囲み
\newtcolorbox{titledbox}[3]{bookbox,colback=#2!6,colframe=#2,boxrule=0pt,
  borderline west={2.5pt}{0pt}{#2},
  title={#3\ #1},fonttitle=\bfseries\sffamily\small,coltitle=#2!85!black,
  attach title to upper={\par\smallskip},colbacktitle=#2!6}
% 章の導入: \begin{chapterlead} ... \end{chapterlead}
\newenvironment{chapterlead}{\begin{titledbox}{この章の位置づけとゴール}{bookblue}{$\blacktriangleright$}}{\end{titledbox}}
% 直感: \begin{intuition}{見出し} ... \end{intuition}
\newenvironment{intuition}[1]{\begin{titledbox}{#1}{bookgreen}{直感}}{\end{titledbox}}
% 注意・よくある誤解: \begin{caution}{見出し} ... \end{caution}
\newenvironment{caution}[1]{\begin{titledbox}{#1}{bookaccent}{注意}}{\end{titledbox}}
% コラム: \begin{column}{見出し} ... \end{column}
\newenvironment{column}[1]{\begin{titledbox}{#1}{bookorange}{コラム}}{\end{titledbox}}
% 実務メモ: \begin{practice}{見出し} ... \end{practice}
\newenvironment{practice}[1]{\begin{titledbox}{#1}{bookpurple}{実務メモ}}{\end{titledbox}}
% 章末のまとめ: \begin{summary} \item ... \end{summary}
\newenvironment{summary}{\section*{この章のまとめ}\addcontentsline{toc}{section}{この章のまとめ}%
  \begin{tcolorbox}[bookbox,colback=figblue!35,borderline west={2.5pt}{0pt}{bookblue}]\begin{itemize}[leftmargin=1.2em]}%
  {\end{itemize}\end{tcolorbox}}
% 演習問題: \begin{exercises} \item ... \end{exercises}
\newenvironment{exercises}{\section*{演習問題}\addcontentsline{toc}{section}{演習問題}%
  \begin{enumerate}[label=\textbf{\thechapter.\arabic*},leftmargin=2.6em]}{\end{enumerate}}
% 付録 D（略解）: \begin{answers}{ch:ml}{章の題} \item ... \end{answers}。番号は章末の演習問題と同じ「章.番号」
\newenvironment{answers}[2]{\section*{第\ref{#1}章\quad #2}%
  \begin{enumerate}[label=\textbf{\ref{#1}.\arabic*},leftmargin=2.6em]}{\end{enumerate}}
% 途中の確認問題: \begin{quiz} ... \end{quiz}
\newenvironment{quiz}{\begin{titledbox}{手を止めて考えてみましょう}{bookgray}{確認}}{\end{titledbox}}

% ---------------------------------------------------------------- 読み方ガイド（初めて読む人のための優先度）
% 優先度の札: \rgmust（必読）, \rgskim（流し読み可）, \rglater（後回し可）
\newcommand{\rgtag}[2]{{\fboxsep=1.3pt\colorbox{#1}{\sffamily\bfseries\scriptsize\textcolor{white}{\hspace{0.15em}#2\hspace{0.15em}}}}}
\newcommand{\rgmust}{\rgtag{bookblue}{必読}}
\newcommand{\rgskim}{\rgtag{bookorange}{流し読み可}}
\newcommand{\rglater}{\rgtag{bookgray}{後回し可}}
\newcommand{\rgpre}{\rgtag{bookgreen}{読む前に}}
% 章の冒頭（chapterlead の直後）に置く道案内:
%   \begin{readingguide}
%   \item[\rgmust] \secref{sec:x}〜\secref{sec:y}：理由。
%   \item[\rgskim] \secref{sec:z}：押さえる結論。
%   \item[\rglater] \secref{sec:w}：戻ってくる目安。
%   \end{readingguide}
\newenvironment{readingguide}{%
  \begin{tcolorbox}[bookbox,lines before break=4,colback=white,colframe=bookblue!70,boxrule=0.5pt,
    title={この章の読み方（初めて読む人へ）},fonttitle=\bfseries\sffamily\small,coltitle=bookblue!85!black,
    attach title to upper={\par\smallskip},colbacktitle=white]\small
  \begin{description}[leftmargin=5.6em,labelwidth=5.0em,labelsep=0.6em,itemsep=1pt,topsep=1pt,parsep=0pt,font=\normalfont]}%
  {\end{description}\end{tcolorbox}}
% 節・小節の見出しの直後に置く 1 行の札（必読の節には何も置かない）:
%   \skimread{結論の「〜」だけ押さえれば先へ進めます。}
%   \laterread{第 14 章で DDP を使うときに戻ってくれば足ります。}
\newcommand{\skimread}[1]{\par\noindent\rgskim\enspace{\small\sffamily\color{black!65}#1}\par\smallskip}
\newcommand{\laterread}[1]{\par\noindent\rglater\enspace{\small\sffamily\color{black!65}#1}\par\smallskip}
% 演習問題の先頭（\begin{exercises} の直後）に置く、初読の人への案内: \exfirst{1 と 3}
\newcommand{\exfirst}[2][]{\item[]\hspace*{-2.6em}\rgtag{bookblue}{初読}\enspace{\small\sffamily\color{black!65}初めて読むときは、#2だけで十分です。#1}}

% ---------------------------------------------------------------- コード
\lstdefinestyle{bookpython}{
  language=Python,
  basicstyle=\ttfamily\footnotesize,
  keywordstyle=\color{bookblue}\bfseries,
  commentstyle=\color{bookgreen!80!black},
  stringstyle=\color{bookaccent},
  showstringspaces=false,
  columns=fullflexible,keepspaces=true,
  breaklines=true,breakatwhitespace=false,
  postbreak=\mbox{\textcolor{bookgray}{$\hookrightarrow$}\space},
  numbers=left,numberstyle=\tiny\color{bookgray},numbersep=6pt,
  xleftmargin=1.2em,
  tabsize=4,upquote=true,
  morekeywords={self,True,False,None,with,as,yield,nonlocal,async,await},
}
% 実行結果・設定ファイル用。Python の色付けをしない。
\lstdefinestyle{booktext}{
  language={},keywordstyle={},commentstyle={},stringstyle={},
  basicstyle=\ttfamily\footnotesize,columns=fullflexible,keepspaces=true,
  breaklines=true,showstringspaces=false,numbers=none,upquote=true,
}
\lstset{style=bookpython}
\newcounter{codelisting}[chapter]
\renewcommand{\thecodelisting}{\thechapter.\arabic{codelisting}}
\newcommand{\snippetdir}{build/snippets}
\newcommand{\codecaption}[2]{%
  \refstepcounter{codelisting}\label{#2}%
  {\small\sffamily\bfseries\color{bookblue}コード\thecodelisting\quad}{\small #1}}
\tcbset{codebox/.style={enhanced,breakable,lines before break=4,sharp corners,colback=codebg,
  colframe=black!25,boxrule=0.4pt,left=1mm,right=1mm,top=0.5mm,bottom=0.5mm,
  before skip=2mm,after skip=3mm}}

% 見出しだけがページ末に残るのを防ぐ（残りが 6 行未満なら改ページ）。
% 先に \penalty0 でページを確定させる（直前の段落があふれた直後は \pagetotal > \pagegoal になるため）。
\newcommand{\codeneedspace}{\par\penalty0\relax\ifdim\pagetotal<\pagegoal
  \ifdim\dimexpr\pagegoal-\pagetotal\relax<6\baselineskip\relax\newpage\fi\fi}
% コード全体を掲載: \codefile{キャプション}{ラベル}{ch10/train.py}
\newcommand{\codefile}[3]{%
  \codeneedspace\medskip\noindent\codecaption{#1}{#2}\hfill{\scriptsize\ttfamily\color{bookgray}code/\detokenize{#3}}\nopagebreak
  \begin{tcolorbox}[codebox]
  \lstinputlisting[style=bookpython]{\snippetdir/#3/_full.txt}
  \end{tcolorbox}}
% コードの一部（スニペット）を掲載: \snippet{キャプション}{ラベル}{ch10/gpt.py}{名前}
\newcommand{\snippet}[4]{%
  \codeneedspace\medskip\noindent\codecaption{#1}{#2}\hfill{\scriptsize\ttfamily\color{bookgray}code/\detokenize{#3}}\nopagebreak
  \begin{tcolorbox}[codebox]
  \lstinputlisting[style=bookpython]{\snippetdir/#3/#4.txt}
  \end{tcolorbox}}
% 実行結果（スクリプトが保存した出力ファイル）: \outputfile{ch10/out/train.txt}
\newcommand{\outputfile}[1]{%
  \begin{tcolorbox}[codebox,colback=outbg,title={実行結果},
    fonttitle=\sffamily\scriptsize\bfseries,coltitle=bookgray,colbacktitle=outbg,
    attach title to upper={\par},top=0mm]
  \lstinputlisting[style=booktext]{code/#1}
  \end{tcolorbox}}
% その場に書く短いコード（説明用の断片・疑似コード）。番号なし。
\lstnewenvironment{pycode}{\lstset{style=bookpython,numbers=none,xleftmargin=0.5em,
  backgroundcolor=\color{codebg},frame=single,rulecolor=\color{black!25}}}{}
% その場に書く実行結果・シェルコマンド・設定ファイルなど。
\lstnewenvironment{textcode}{\lstset{style=booktext,xleftmargin=0.5em,
  backgroundcolor=\color{outbg},frame=single,rulecolor=\color{black!25}}}{}

% ---------------------------------------------------------------- TikZ 共通スタイル
\tikzset{
  >={Stealth[length=2mm]},
  box/.style={draw=bookblue,fill=figblue,rounded corners=2pt,minimum height=7mm,
    minimum width=16mm,align=center,font=\small},
  redbox/.style={box,draw=bookaccent,fill=figred},
  greenbox/.style={box,draw=bookgreen,fill=figgreen},
  orangebox/.style={box,draw=bookorange,fill=figorange},
  purplebox/.style={box,draw=bookpurple,fill=figpurple},
  graybox/.style={box,draw=black!45,fill=figgray},
  plainbox/.style={box,draw=black!60,fill=white},
  op/.style={draw=black!70,circle,inner sep=1pt,minimum size=5mm,font=\small,fill=white},
  arr/.style={->,thick,draw=black!70},
  darr/.style={->,thick,dashed,draw=black!55},
  lbl/.style={font=\footnotesize,text=bookgray},
  neuron/.style={draw=bookblue,fill=figblue,circle,minimum size=6mm,inner sep=0pt},
  tok/.style={draw=black!50,fill=figorange,rounded corners=1pt,minimum height=6mm,
    minimum width=9mm,font=\small,inner sep=2pt},
  cell/.style={draw=black!40,minimum size=6mm,font=\footnotesize,inner sep=0pt},
}
\pgfplotsset{bookplot/.style={width=0.78\linewidth,height=0.48\linewidth,
  grid=major,grid style={black!10},tick label style={font=\footnotesize},
  label style={font=\small},legend style={font=\footnotesize,draw=black!30},legend cell align=left,
  every axis plot/.append style={thick}},
  every axis/.append style={bookplot}}

% ---------------------------------------------------------------- 全体地図
% \bookmap{k}: 本書の 7 部構成を帯で示し、第 k 部を強調する。各章冒頭に置く。
\newcommand{\bookmapitem}[4]{% #1 番号 #2 現在の部 #3 見出し #4 x 位置
  \ifnum#1=#2
    \node[draw=bookblue,fill=bookblue,text=white,rounded corners=1.5pt,
      minimum height=9mm,text width=15.5mm,align=center,font=\sffamily\bfseries\tiny,inner sep=1pt]
      (bm#1) at (#4,0) {#3};
  \else
    \node[draw=black!30,fill=black!4,text=black!60,rounded corners=1.5pt,
      minimum height=9mm,text width=15.5mm,align=center,font=\sffamily\tiny,inner sep=1pt]
      (bm#1) at (#4,0) {#3};
  \fi}
\newcommand{\bookmap}[1]{%
  \begin{center}\begin{tikzpicture}[x=1mm,y=1mm]
    \bookmapitem{1}{#1}{I\ 準備\\数学と道具}{0}
    \bookmapitem{2}{#1}{II\ 機械学習と\\ニューラルネット}{18.5}
    \bookmapitem{3}{#1}{III\ 言語モデルと\\Transformer}{37}
    \bookmapitem{4}{#1}{IV\ 大規模学習\\GPU・分散}{55.5}
    \bookmapitem{5}{#1}{V\ 事後学習\\微調整・蒸留}{74}
    \bookmapitem{6}{#1}{VI\ 強化学習}{92.5}
    \bookmapitem{7}{#1}{VII\ エージェント\\と総仕上げ}{111}
    \foreach \a/\b in {1/2,2/3,3/4,4/5,5/6,6/7}
      \draw[->,black!45,line width=0.4pt] (bm\a) -- (bm\b);
  \end{tikzpicture}\end{center}}
